\scanpage{12}
Soit maintenant $\Pi$ hyperpolyèdre comb. de dim $n$, $G$ le groupe de ses autom, $R_\Pi$ $G_{\deleted{\unsure{g|G}}}$-torseur des repères -- son commutant est un quotient $\unsure{\mathscr G_\Pi|G_\Pi|\mathfrak G_\Pi}$ de $\mathscr G$,

$C_{pr}(\Pi)$ cat des réa. géom régulières \inserted{proj} non dégénérées de $\Pi$ ($\Rightarrow\varphi_G:G\longrightarrow\operatorname{Aut}(X)$).
\deleted{$\unsure{C_{pr}(\Pi)}$}

A' tt $r\in R$ correspond un objet de $C_{pr}$ (cas universel épinglé) et pour $r\in R$ variable, on a un système transitif d'iso. entre ces objets. Donc $C_{pr}(\Pi)$ s'identifie \inserted{(pour $\Pi$ fixé)} à \deleted{la cat.} \deleted{des couples \unsure{$(\ )$|$(\ldots)$}} \inserted{sous cat. pleine} de $C_{pr}$ formé des $(X,\varphi_{\mathscr G},\ldots)$ tels que $\varphi_{\mathscr G}:\mathscr G\longrightarrow\operatorname{Aut}(X)$ se factorise par le quotient $\mathscr G_\Pi$ de $\mathscr G$. On fait opérer à \uline{droite} $\mathscr G_\Pi$ sur $R_\Pi$ via son action à gauche, et on l'utilise pour tordre $(X,\varphi_{\mathscr G}\ldots)$ en
\[
R_\Pi\mathbin{\wedge_{\mathscr G_\Pi}} X
\]

Même procédé pour les autres descriptions -- pour $C'_{pr}(\Pi)$ etc ....

Autre point de vue : On ne se donne pas $\Pi$, mais seulement un \inserted{groupe} quotient déterminé $\mathscr G_0$ de $\mathscr G$, \deleted{\unsure{les}} correspondant à un hyperpolyèdre régulier combinatoire épinglé. On considère un hyperpolyèdre régulier

% archive sheet number: 11

